\section{The Autonomous Drift Problem}

\subsection{Definitions}

In a recursively spawned agent architecture, a \emph{chain of agency} is a directed graph
$G = (V, E)$ where each vertex $v_i \in V$ represents an agent instance and each directed
edge $(v_i, v_j) \in E$ denotes a spawn relationship: $v_i$ created $v_j$.
Every agent $v_i \neq v_{\text{root}}$ has exactly one incoming spawn edge;
the root agent $v_0$ has none.

\begin{definition}[Orphaned Agent]
An agent $v$ is \emph{orphaned} if its parent $p(v)$ has either (a) terminated
normally, (b) crashed, or (c) become unreachable due to network partition or timeout.
Orphaning severs the chain's upward reference without terminating $v$ itself.
\end{definition}

\begin{definition}[Human Anchor]
A chain has a \emph{human anchor} if there exists at least one agent $a \in V$
such that $a$ holds a confirmed, responsive communication channel to a human
principal. Let $H \subseteq V$ denote the set of such agents.
\end{definition}

\begin{definition}[Drifted Agent]
An agent $v$ is \emph{drifted} if the chain containing $v$ has no path from any
agent in the connected component of $v$ to any agent in $H$. Formally, letting
$C(v)$ be the connected component of $v$ in $G$:
\[
\text{drifted}(v) \iff C(v) \cap H = \emptyset .
\]
\end{definition}

\begin{definition}[Drift Event]
A \emph{drift event} occurs at time $t$ when an agent $v$ transitions from
$\neg\text{drifted}(v)$ to $\text{drifted}(v)$.
\end{definition}

A chain may be long, shallow, or branched. Drift is independent of chain length
or branching factor; it is determined solely by the existence of a path to $H$.

\subsection{Conditions That Cause Drift}

Drift is not a single failure mode but the outcome of several distinct conditions
that interrupt the chain's continuity.

\medskip
\noindent\textbf{Parent Termination.}
The most direct cause: an agent $p$ spawns child $c$ and then exits. If $p$ is the
last anchor-bearing ancestor in $c$'s component, $c$ becomes drifted upon $p$'s
exit. This is the common case in any architecture that treats spawn as ``fire and
forget.''

\medskip
\noindent\textbf{Communication Partition.}
Even when the parent remains alive, prolonged network latency or transient partition
can make the child unable to confirm liveness or route messages to the parent.
If the child's refresh interval exceeds its timeout threshold, the child must
assume the parent is unavailable and may enter a degraded or drifted state.

\medskip
\noindent\textbf{Cascading Orphaning.}
An agent may spawn multiple children. If the parent terminates or partitions,
all children become simultaneously orphaned. If none of those children have
independent channels to $H$, all drift. This is a burst event, not a sequential one.

\medskip
\noindent\textbf{Ambiguous Liveness.}
An agent whose parent crashes without broadcasting a termination signal enters a
state where the child cannot determine whether the parent is dead or merely slow.
Conservative protocols may treat this as a partition and trigger drift countermeasures;
optimistic protocols may wait indefinitely, during which the child is effectively
drifted from the perspective of any human principal expecting a response.

\subsection{Failure Modes}

Once an agent is drifted, the failure modes that follow are both predictable and
escalating.

\medskip
\noindent\textbf{Unbounded Work Continuation.}
A drifted agent that does not detect its drift continues executing its task
indefinitely. Unlike a supervised process that crashes or halts, it neither
produces output accessible to any human nor responds to any principal.
In the worst case, it consumes resources (compute, tokens, API quota) with
no possible return.

\medskip
\noindent\textbf{Conflicting Outputs.}
A drifted agent may produce artifacts, database writes, or API calls that
conflict with parallel or subsequent work in the parent chain. Because the
drifted agent cannot receive correction signals from any anchor, conflicts
accumulate without resolution.

\medskip
\noindent\textbf{Lost Lineage.}
When a drifted agent terminates (normally or abnormally), all records of its
spawn tree are lost if the agent did not proactively maintain a durable lineage
log. The system can no longer reconstruct which agents existed, what they did,
or why they were spawned.

\medskip
\noindent\textbf{Security Surface.}
A drifted agent retains whatever permissions its parent chain possessed. An agent
that drifted due to a partition may still hold valid credentials, API tokens,
or file system access that it exercises without oversight. This creates an
unmonitored privilege window.

\subsection{Why Depth Limits Alone Do Not Solve Drift}

A common architectural response to unbounded spawning is to impose a depth limit
$D_{\max}$ on the agent chain. The intuition is that if depth is bounded,
drift is contained. This intuition is incorrect for three independent reasons.

\medskip
\noindent\textbf{Drift is a function of connectivity, not length.}
An agent at depth 2 with a live channel to its human anchor is not drifted.
An agent at depth 1 with no such channel \emph{is} drifted. Depth limits
constrain the size of the chain, not its attachment to an anchor. A shallow
chain with zero anchor connections is fully drifted; a deep chain in which
every level maintains a channel to $H$ is never drifted.

Formally, drift$(v) \iff C(v) \cap H = \emptyset$. Depth does not appear in
this condition. Therefore a depth limit $D_{\max}$ does not change the set of
drifted agents unless $D_{\max}$ is set to zero.

\medskip
\noindent\textbf{Width is a free variable.}
A depth-limited architecture can still spawn unbounded width at each level.
An agent at depth $D_{\max} - 1$ can spawn arbitrarily many children at depth
$D_{\max}$. If the parent at depth $D_{\max} - 1$ is not connected to $H$,
all of its children are drifted regardless of the depth limit. The bound on
depth does not limit the blast radius of a single orphaning event.

\medskip
\noindent\textbf{Depth limits do not preserve anchors.}
Consider a chain of depth $D_{\max}$ where the root agent is the sole anchor.
If the root terminates, all $D_{\max}$ levels become drifted simultaneously.
Imposing a depth limit of $D_{\max} / 2$ merely delays the onset; the drift
event occurs when the last anchor-bearing ancestor exits, which is independent
of how many levels exist below it.

The only way a depth limit meaningfully reduces drift risk is if each depth
level maintains an independent channel to $H$. In that case, the depth limit
acts as a \emph{replication} mechanism, not a containment mechanism. This is
substantially different from a simple depth cap and requires explicit anchor
redundancy design.

\subsection{A Simple Formal Model}

We formalize the above to make the dependencies precise.

Let $P$ be the set of all possible agent processes. Each process $p \in P$ has:

\begin{itemize}
    \item $\text{parent}(p) \in P \cup \{\bot\}$ --- the spawning agent, or $\bot$ for the root
    \item $\text{alive}(p, t) \in \{\text{true}, \text{false}\}$ --- liveness at time $t$
    \item $\text{connected}(p, H, t) \in \{\text{true}, \text{false}\}$ --- whether $p$ has a live channel to any human anchor
    \item $\text{drifted}(p, t) \in \{\text{true}, \text{false}\}$ --- drift state at $t$
\end{itemize}

The drift state is defined recursively:

\begin{equation}
\text{drifted}(p, t) =
\begin{cases}
\text{false} & \text{if } \text{connected}(p, H, t) = \text{true} \\[6pt]
\text{true}  & \text{if } \text{alive}(\text{parent}(p), t) = \text{false} \text{ and } \text{connected}(p, H, t) = \text{false} \\[6pt]
\text{drifted}(\text{parent}(p), t) \land \neg\text{connected}(p, H, t) & \text{otherwise}
\end{cases}
\end{equation}

That is: an agent is drifted if it has no anchor connection \emph{and} its
parent is either dead or itself drifted. This is a recursive definition that
traces upward until either a living anchor-connected ancestor is found or the
root is reached.

\medskip
\noindent\textbf{Drift propagation theorem.}
\textit{If an agent $a$ becomes drifted at time $t$, then every descendant
$d$ of $a$ (in the spawn tree) for which $\text{connected}(d, H, t) = \text{false}$
also becomes drifted at time $t$.}

\textit{Proof.} By induction on the depth distance from $a$ to $d$. For the
direct child $c$ of $a$, $\text{parent}(c) = a$ and $a$ is drifted, so the third
case of the drift definition applies: $\text{drifted}(c, t) = \text{true}$.
Inductive step follows identically. $\square$

This theorem is the core reason depth limits fail: a single drift event at any
point in the chain propagates downward to all unattached descendants,
regardless of how deep the tree is.

\medskip
\noindent\textbf{Correction condition.}
An agent $p$ can be \emph{re-anchored} at time $t > t_{\text{drift}}$ if either:
\begin{enumerate}
    \item A human principal establishes a new channel directly to $p$ (manual intervention), or
    \item An ancestor $a$ of $p$ is re-anchored and re-establishes connectivity to $p$.
\end{enumerate}

This defines the only two recovery paths from drift. Neither happens automatically.
Automatic drift detection without automatic re-anchoring is an incomplete solution.

\end{document}
